\section{Segurança e controle de acesso}\label{sec:seguranca}

A segurança foi considerada como um requisito transversal do sistema, especialmente em razão da natureza das informações armazenadas durante o processo de avaliação psicopedagógica. A aplicação possui diferentes perfis de usuário e mecanismos de autenticação e autorização, buscando restringir o acesso às funcionalidades de acordo com as responsabilidades atribuídas a cada usuário.

A autenticação da aplicação foi implementada por meio de JSON Web Tokens (JWT), utilizando tokens de acesso e tokens de atualização (refresh tokens). Após a autenticação, o usuário recebe as credenciais necessárias para acessar os recursos protegidos da API. O backend utiliza middleware de autenticação para verificar as requisições direcionadas às rotas que exigem usuário autenticado, impedindo o acesso direto a esses recursos sem as credenciais correspondentes.

Além da autenticação, o sistema possui controle de acesso baseado em perfis. Foram definidos os papéis \textit{ADMIN} e \textit{PSICOPEDAGOGO}, permitindo diferenciar as permissões disponíveis para cada tipo de usuário. O perfil administrador possui acesso às funcionalidades de gerenciamento de usuários e a todos os registros do sistema, enquanto o perfil de psicopedagogo pode realizar operações relacionadas aos aprendentes, avaliações, aplicações e relatórios sob sua responsabilidade.

O controle por perfil é aplicado no backend por meio do middleware \textit{authorize}, presente nas rotas que exigem permissões específicas. Dessa forma, tentativas de acessar o gerenciamento de usuários ou de excluir instrumentos avaliativos, operações exclusivas do perfil \textit{ADMIN}, são bloqueadas com resposta 403 quando realizadas por um usuário com perfil \textit{PSICOPEDAGOGO}, antes de chegar à camada de serviço. 

A interface complementa esse controle ao condicionar a exibição de determinados elementos ao perfil do usuário, como o menu de gerenciamento de usuários e o botão de exclusão de instrumentos, exibidos somente ao administrador. O controle definitivo, entretanto, é exercido pelo backend, de modo que a restrição na interface funciona como auxílio à experiência do usuário, não como mecanismo de segurança principal.

Além do controle por perfil, o sistema implementa restrição de visibilidade baseada em titularidade. Aprendentes e avaliações são vinculados ao usuário que os criou por meio dos campos \textit{responsavelId} e \textit{autorId}, respectivamente, e as aplicações e os relatórios seguem a titularidade do aprendente a que pertencem. Nas operações de listagem, consulta, alteração e exclusão, a camada de serviço do backend aplica automaticamente um filtro que restringe os registros ao usuário autenticado, e a criação de uma aplicação exige que o aprendente e o instrumento pertençam ao próprio profissional. 

Registros de outros profissionais são tratados como inexistentes, sem revelar que existem. O perfil \textit{ADMIN} constitui exceção a essa regra, tendo acesso irrestrito a todos os registros do sistema. Esse mecanismo garante que profissionais distintos não acessem os registros uns dos outros, mesmo possuindo o mesmo perfil de acesso.

Essa diferenciação também está representada no banco de dados, na entidade \textit{users}, que possui o campo \textit{role} para identificação do perfil do usuário e o campo \textit{status}, responsável por indicar se a conta está ativa ou inativa. Dessa maneira, o controle de acesso considera tanto a identidade do usuário quanto as permissões associadas ao seu perfil.

No frontend, o controle de acesso é complementado pelo uso do \textit{AuthContext} e do componente \textit{ProtectedRoute}. O contexto de autenticação mantém as informações necessárias para o gerenciamento da sessão do usuário, enquanto as rotas protegidas restringem o acesso às páginas que dependem de autenticação. A interface também considera os diferentes perfis de acesso para disponibilizar as funcionalidades correspondentes a cada usuário.

A comunicação entre frontend e backend é realizada por meio da API REST, sendo as operações protegidas verificadas pelo backend. Essa separação permite que o controle de segurança não dependa exclusivamente das restrições apresentadas na interface, mantendo a validação das permissões no lado servidor.

O gerenciamento das credenciais também conta com estruturas específicas para tokens de atualização e recuperação de senha. A entidade \textit{refresh\_tokens} armazena os tokens associados aos usuários, juntamente com seus prazos de validade e informações de revogação. Já a entidade \textit{password\_resets} permite registrar tokens destinados ao processo de recuperação de senha, incluindo prazo de validade e indicação de utilização.

Outro mecanismo utilizado na implementação é a validação dos dados recebidos pela API. O backend utiliza Zod para validação das entradas, permitindo verificar se os dados enviados às operações atendem às estruturas esperadas antes de serem processados. Esse mecanismo complementa as medidas de autenticação e autorização ao contribuir para o controle das informações recebidas pela aplicação.

Para a rastreabilidade, os campos de controle temporal (\textit{createdAt} e \textit{updatedAt}) e os vínculos de autoria do modelo (\textit{responsavelId}, \textit{autorId} e \textit{aplicadorId}) não são suficientes quando se trata de dados de saúde de crianças. Eles indicam quem criou cada registro e quando ele foi alterado pela última vez, mas não registram quem fez cada alteração, quais eram os valores anteriores nem quem consultou as informações. Este último ponto é justamente o que a Lei Geral de Proteção de Dados e o sigilo profissional mais exigem. Além disso, o perfil \textit{ADMIN} tem acesso irrestrito a todos os dados clínicos, e esse acesso não deixaria rastro.

Por esses motivos, o sistema mantém um registro de auditoria na entidade \textit{audit\_logs}, preenchido automaticamente por uma extensão do Prisma Client, conforme descrito na Seção \ref{sec:backend}. O registro atende aos quatro pontos acima:

\begin{itemize}
    \item cada criação, alteração e exclusão identifica o usuário que a realizou, o seu perfil, o IP de origem e o momento da operação;
    \item as alterações guardam os valores anteriores e os novos dos campos modificados, e as exclusões guardam o conteúdo excluído;
    \item as consultas e listagens de dados clínicos (aprendentes, histórico de acompanhamento, aplicações e relatórios) são registradas e associadas ao aprendente, permitindo saber quem acessou os dados de cada criança;
    \item os acessos do administrador são registrados como os de qualquer outro usuário, e o psicopedagogo responsável pode consultá-los na aba ``Registro de Acessos'' do aprendente.
\end{itemize}

Os logins bem-sucedidos e as tentativas recusadas também são registrados, e o hash da senha nunca é gravado no registro.

O registro de auditoria, entretanto, possui limitações que precisam ser consideradas:

\begin{itemize}
    \item o registro é gravado após a operação, em uma conexão separada, e não faz parte da mesma transação. Se a gravação da auditoria falhar, a operação é mantida e a falha aparece apenas no log do servidor. Em operações executadas dentro de uma transação, o registro pode permanecer mesmo que a transação seja desfeita;
    \item não há garantia de imutabilidade. A aplicação não oferece meios de alterar ou excluir o registro, mas quem tem acesso direto ao banco de dados pode modificá-lo, pois não há encadeamento criptográfico dos registros nem armazenamento externo somente de gravação;
    \item acessos feitos diretamente ao banco de dados ou às suas cópias de segurança não passam pela aplicação e, portanto, não são registrados;
    \item o próprio registro contém dados sensíveis, pois guarda valores anteriores e novos dos dados clínicos. Por isso, a consulta geral é restrita ao administrador, mas não foi definida uma política de retenção e descarte desses registros;
    \item o administrador continua com acesso irrestrito: a auditoria torna esse acesso rastreável e visível ao profissional responsável, mas não o impede nem exige justificativa. As consultas do administrador ao próprio registro de auditoria não são registradas;
    \item as leituras feitas internamente pelo sistema, como a verificação de permissão antes de uma alteração, também são registradas como consultas, e as listagens geram um registro por item exibido. Isso aumenta o volume de registros e pode dificultar a análise. Contagens agregadas, como as do dashboard, não são registradas, assim como as tentativas de acesso negadas a registros de outros profissionais;
    \item o endereço IP registrado depende da configuração do proxy reverso. Se a API estiver acessível diretamente, sem passar pelo proxy, o cabeçalho que informa o IP de origem pode ser falsificado.
\end{itemize}

Essas limitações indicam que o registro implementado constitui um mecanismo básico de rastreabilidade, adequado ao estágio de protótipo do sistema, mas não substitui as medidas organizacionais e técnicas exigidas para o uso com dados reais, como a imutabilidade do registro, a política de retenção e o controle do acesso direto à infraestrutura.

\subsection{Adequação à LGPD e limitações para uso em produção}

O sistema armazena dois tipos de dados com proteção legal reforçada pela Lei Geral de Proteção de Dados. Os dados dos aprendentes envolvem crianças, categoria sujeita ao art. 14 da LGPD, que exige tratamento no melhor interesse do menor e, em regra, consentimento do responsável legal. As respostas das avaliações e os relatórios podem conter dados de saúde, classificados como dados pessoais sensíveis pelo art. 5º, inciso II, e sujeitos às restrições do art. 11, que limita as bases legais admissíveis para seu tratamento.

Do ponto de vista da minimização de dados, os campos opcionais da entidade \textit{aprendentes} — contato secundário, e-mail, escola, série e turno — permitem que o profissional registre apenas as informações necessárias ao atendimento. Os campos obrigatórios (nome, data de nascimento, sexo, responsável e contato principal) foram definidos com base nas informações mínimas identificadas como necessárias ao processo avaliativo. Ainda assim, a decisão sobre quais campos preencher é do profissional, e o sistema não impede o registro de informações além do estritamente necessário.

As seguintes medidas foram implementadas em relação à proteção dos dados: controle de acesso por perfil e titularidade, registro de auditoria das operações e acessos, e exclusão lógica dos registros. Contudo, o sistema apresenta limitações relevantes para uso com dados reais que precisam ser reconhecidas:

\begin{itemize}
    \item \textbf{Consentimento do responsável legal}: o sistema não possui mecanismo para registrar o consentimento do responsável da criança para o tratamento de seus dados, exigido pelo art. 14 da LGPD. Essa responsabilidade recai sobre o profissional que utiliza o sistema, devendo ser gerenciada por procedimentos externos à plataforma;
    \item \textbf{Base legal}: o sistema não registra a base legal utilizada para o tratamento dos dados de cada aprendente. Para dados de saúde de crianças, as bases admissíveis são restritas (art. 11, LGPD), e sua documentação é necessária para eventual prestação de contas (\textit{accountability});
    \item \textbf{Criptografia em repouso}: os dados armazenados no banco de dados não são criptografados em nível de campo. A proteção depende integralmente dos controles de acesso à infraestrutura (autenticação no banco, permissões do sistema operacional). Para uso com dados reais, a criptografia em repouso dos campos sensíveis — especialmente os campos de texto livre de aprendentes, históricos, respostas e relatórios — constitui medida adicional recomendada;
    \item \textbf{Backups}: não foi definida uma política de backup. Cópias de segurança do banco de dados contêm os mesmos dados sensíveis que a aplicação principal e devem ser protegidas com os mesmos controles de acesso. Acessos a backups não passam pela aplicação e, portanto, não são registrados na auditoria;
    \item \textbf{Retenção e descarte}: não foi definida uma política de retenção e descarte dos dados. A LGPD exige que os dados sejam eliminados ao término do tratamento ou mediante solicitação do titular. O sistema implementa exclusão lógica (campo \textit{deletedAt}), mas não possui mecanismo de exclusão definitiva dos dados nem de seus registros de auditoria.
\end{itemize}

Essas limitações são inerentes ao estágio de protótipo do sistema e não inviabilizam sua avaliação funcional. Para uso em produção com dados reais de crianças, seria necessário endereçar cada um desses pontos, além de elaborar um Relatório de Impacto à Proteção de Dados Pessoais (RIPD) conforme previsto no art. 38 da LGPD.

Assim, a segurança implementada combina autenticação por JWT, controle de acesso por perfis, restrição de visibilidade por titularidade, proteção das rotas, gerenciamento de sessão, validação das informações recebidas pela API e registro de auditoria. Esses mecanismos estabelecem uma camada de proteção adequada ao estágio de protótipo do sistema, mas não substituem as medidas organizacionais e técnicas exigidas para o tratamento de dados pessoais sensíveis de crianças em ambiente de produção.